\section{A 639-bit equational packet using a descriptor as its own class}
\label{sec:e639}

A separate result improves the complete equational input from the pinned
663-bit packet to \textbf{639 bits}.  The seven equations, nine operation
symbols, four local variables, and ground query are all charged under the
unchanged codec.  The change removes six equation-term nodes by eliminating
three uses of the descriptor projection $V(b,d)$.  The descriptor itself
now serves as the represented class.

The mathematical issue is not the textual replacement.  A descriptor must
carry both an arbitrary formula and arbitrary parameters while preserving
the formula's interpretation.  We prove that a duplicated parameter spine
can carry those parameters with constant class truth value.  This gives a
finite, noncircular construction of every required descriptor.  We also
recheck the upper-model construction: its raw-value choices still depend
only on proper subtrees.  The conclusion is a new relative calibration,
not a claim that the old and new presentations have identical models.

No binary Turing-machine state bound follows from 639 bits.  The result
provides a smaller complete input for a future equational search engine;
that engine, its enumeration, and its blank-tape initialization remain to
be implemented and counted.

\subsection{Literal signature, axioms, and query}

The ordered signature is
\[
 d^0,N^2,V^2,P^2,A^2,Q^1,W^2,p_E^2,p_\varepsilon^2.
\]
All equations below are universally quantified in their displayed
variables.  The first six are unchanged from the E663 packet:
\begin{align}
 N(N(x,N(N(y,x),x)),N(y,N(z,x)))&=y,\label{eq:e639-sh1}\\
 V(P(x,s),d)&=x,\label{eq:e639-head}\\
 V(P(x,s),Q(j))&=V(s,j),\label{eq:e639-tail}\\
 p_\varepsilon(s,P(a,b))
   &=N(p_\varepsilon(s,a),p_\varepsilon(s,b)),\label{eq:e639-nand}\\
 p_\varepsilon(s,A(i,j))
   &=p_E(V(s,i),V(s,j)),\label{eq:e639-atom}\\
 p_\varepsilon(s,Q(a))
   &=N\bigl(p_\varepsilon(P(W(s,a),s),a),
              p_\varepsilon(P(z,s),a)\bigr).\label{eq:e639-binder}
\end{align}
The first equation is the published Sh1 single axiom for Boolean algebra,
due to McCune, Veroff, Fitelson, Harris, Feist, and Wos; it is inherited,
not a newly discovered Boolean axiom.\footnote{Primary paper:
\url{https://www.cs.unm.edu/~mccune/papers/basax/v12.pdf},
\emph{Short Single Axioms for Boolean Algebra}, Journal of Automated
Reasoning 29 (2002), 1--16.}
Thus $N$ is NAND in a Boolean algebra on the whole carrier.  Its top
and bottom are written $1$ and $0$.  The distinguished native element
$d$ is not assumed to equal $1$ until that equality is proved below.

The seventh equation has the same Boolean kernel as E663, but uses the
new direct descriptor relation
\begin{equation}
 \widetilde R_b(x,y)=p_\varepsilon(P(y,P(x,b)),b).
 \label{eq:e639-direct}
\end{equation}
The following are abbreviations only; the bit packet expands them all:
\begin{align*}
 h&=Q(b),& n_z&=P(b,z),&m_z&=A(b,z),\\
 r&=\widetilde R_b(h,n_z),&
 e&=p_E(m_z,n_z),&f&=p_E(m_z,z),\\
 a&=p_E(W(x,d),x),&b_0&=p_E(x,h),&c&=\widetilde R_b(x,y),\\
 u&=p_E(y,h),&v&=\widetilde R_b(P(b,y),y).
\end{align*}
Define
\begin{align*}
 K&=N(c,N(a,N(d,b_0))),\\
 H&=N\bigl(N(N(d,u),K),N(u,v)\bigr),\\
 B^-&=N\bigl(r,N(d,N(e,N(d,f)))\bigr),\\
 F(d,r,e,f,a,b_0,c,u,v)&=N(B^-,N(d,H)).
\end{align*}
The last equation is
\begin{equation}
                d=F(d,r,e,f,a,b_0,c,u,v).
 \label{eq:e639-closure}
\end{equation}
The designated ground query remains
\begin{equation}
                         Q(d)=d.
 \label{eq:e639-query}
\end{equation}

An equational theory always has a singleton model.  Accordingly the
consistency question here is noncollapse: equivalently, the consistency
of the ordinary first-order theory consisting of these seven equations
and the ground disequality $Q(d)\ne d$.  It is not the consistency of
the equations alone.

\subsection{A class-transparent parameter spine}

For a native element $a$ and environment $s$, write
$\varepsilon_s(a)=p_\varepsilon(s,a)$, with its value in the full
Boolean reduct.  All identities in this subsection hold before choosing
an ultrafilter and without assuming any value for $\varepsilon_s(d)$.
Put
\[
 T=P(d,P(d,d)).
\]
For any finite tuple of native parameters $p_1,\ldots,p_k$, define
\[
 S_{k+1}=T,\qquad
 S_i=P(p_i,P(p_i,S_{i+1}))\quad(1\leq i\leq k).
\]
In the empty-parameter case let $S_1=T$.

\paragraph{Masking lemma.}
For every environment $s$ and every $i$,
\[
 \varepsilon_s(T)=1,\qquad \varepsilon_s(S_i)=1.
\]
For every additional native class $c$, the descriptor
\begin{equation}
                 b=P(P(c,c),S_1)
 \label{eq:e639-maskdescriptor}
\end{equation}
satisfies
\begin{equation}
                         \varepsilon_s(b)=\varepsilon_s(c)
 \label{eq:e639-transparency}
\end{equation}
for every $s$.

\paragraph{Proof.}
Equation~\eqref{eq:e639-nand} and Boolean algebra give
\[
 \varepsilon_s(T)
 =N(t,N(t,t))=1,\qquad t=\varepsilon_s(d).
\]
If $\varepsilon_s(S_{i+1})=1$, then
\[
 \varepsilon_s(S_i)=N(q,N(q,1))=N(q,\neg q)=1,
 \qquad q=\varepsilon_s(p_i).
\]
Finite descending induction proves the spine claim.  Finally,
\[
 \varepsilon_s(b)=N(N(q,q),1)=q,
 \qquad q=\varepsilon_s(c).
\]
These are Boolean identities applied to already evaluated predicate
values; they impose no compatibility of $P$ or the other native
functions with a Boolean homomorphism.\hfill$\square$

Set $\nu_0=d$ and $\nu_{j+1}=Q(\nu_j)$.  Repeated use of the
native tail equation followed by the head equation shows that
$V(s,\nu_j)$ reads coordinate $j$ of any sufficiently long finite
$P$-stack, with the head numbered zero.  For
$s=P(y,P(x,b))$ and the descriptor~\eqref{eq:e639-maskdescriptor},
the relevant coordinates are exactly
\[
\begin{array}{c|ccccccc}
\text{coordinate}&0&1&2&3&4&5&6\\ \hline
\text{value}&y&x&P(c,c)&p_1&p_1&p_2&p_2.
\end{array}
\]
Generally $p_i$ occurs at coordinates $2i+1$ and $2i+2$.
No assertion about coordinates beyond the finite stack is needed.

\paragraph{Representation lemma.}
After passage to two truth values, every finite equality-free
pure-membership formula $R(x,y;p_1,\ldots,p_k)$ is represented
by $\widetilde R_b$ for a descriptor of the form
\eqref{eq:e639-maskdescriptor}.

\paragraph{Proof.}
Build its finite class code $c$ first.  Atomic membership between
coordinates $i,j$ is represented by $A(\nu_i,\nu_j)$, NAND
by $P$, and ordinary existential quantification by $Q(P(a,a))$.
The binder equation provides existential negation, as checked below.
Use coordinate zero for the output $y$, coordinate one for the input
$x$, and coordinate $2i+1$ for parameter $p_i$.  Coordinate two
and the duplicate parameter coordinates are unused.  Under a binder,
shift all existing coordinate indices by one and use the new head
coordinate for the bound variable.  Structural induction on the
finite formula proves this representation.  Boolean truth and falsity
can be expressed by NAND tautologies and contradictions if the
formula language includes them.

Now form $S_1$ and $b$ using the already built code $c$.  The
coordinate identities put exactly the desired parameters in
$P(y,P(x,b))$.  Identity~\eqref{eq:e639-transparency} replaces
evaluation of class $b$ by evaluation of $c$ at this same environment.
Thus $\widetilde R_b(x,y)$ is $R(x,y;\bar p)$.
There is no circular definition: $c$ only mentions fixed finite
coordinate indices and is built before $b$.  The arbitrary class
truth values of the native parameters are canceled by the spine.
\hfill$\square$

\subsection{Truth forcing and extraction of raw EPSC}

The inherited Boolean kernel has the cofactor identity
\begin{equation}
 \neg d\wedge F(d,r,e,f,a,b_0,c,u,v)=\neg d\wedge r.
 \label{eq:e639-cofactor}
\end{equation}
Indeed, at $d=0$ every $N(d,t)$ is $1$, so $B^-=\neg r$
and $F=r$.  Projecting evaluated Boolean values to the principal
interval $[0,\neg d]$ proves the identity for an arbitrary Boolean
algebra.  If $r=1$ and $d=F$, it follows that $\neg d=0$.

In the new equation we obtain $r=1$ simply by substituting $b=T$:
the masking lemma gives $p_\varepsilon(s,T)=1$ for every $s$,
including the environment in the displayed $r$.  Therefore
every model of the seven equations satisfies $d=1$.
This truth-forcing argument no longer uses the descriptor-head
projection; its changed ground instance is explicitly supplied in
the artifact.

For completeness we state the raw theory to which the packet is
calibrated.  In a nonempty structure $(M,E)$ with no native equality
in the formula schema, put
\begin{align*}
 \mathsf{Empty}(x)&\iff\forall t\,\neg E(t,x),\\
 y\subseteq_E z&\iff\forall t\,(E(t,y)\to E(t,z)),\\
 \mathsf{LB}(R)&\iff\forall x\,\exists z\,\forall y\,
                          (R(x,y)\to y\subseteq_E z).
\end{align*}
For each actual finite first-order matrix $R(x,y;\bar p)$, raw
empty-input positive supported closure, abbreviated EPSC, is
\begin{equation}
\mathsf{LB}(R)\ \longrightarrow\
\exists h\left[
\begin{aligned}
 &\forall x,y\,
  (((\mathsf{Empty}(x)\lor E(x,h))\land R(x,y))\to E(y,h))\\
 &{}\land\forall y\,(E(y,h)\to\exists x\,R(x,y))
\end{aligned}\right].
\label{eq:e639-epsc}
\end{equation}
Parameters are allowed.  This quantifies over syntactic formula
instances, not all external relations on the domain.

\paragraph{Extraction theorem.}
Every nontrivial model of the seven new equations yields a raw-EPSC
structure on the same carrier.

\paragraph{Proof.}
Choose a Boolean homomorphism from the nontrivial Boolean reduct
onto the two-element algebra.  Apply it to the values of $p_E$
and $p_\varepsilon$, while retaining the native operations on
the original carrier with their actual equality.

For the binder, let $w$ be the truth value at the selected witness
$W(s,a)$, let $u_z$ be the value at the arbitrary $z$, and let $q$
be the value of $p_\varepsilon(s,Q(a))$.  The binder equation says
$q=\neg(w\land u_z)$ for every $z$.  Substituting the selected
witness gives $q=\neg w$, and the remaining equations give
$w\leq u_z$.  Hence $q$ is true exactly when some body instance
is false.  This proves the existential-negation interpretation of
$Q$ used in the representation lemma.

The derived equation $d=1$ makes the last equation precisely
\[
 (r\land e\land\neg f)\ \lor\
 \bigl((a\lor\neg c\lor u)
       \land(\neg b_0\lor\neg c\lor u)
       \land(\neg u\lor v)\bigr).
\]
Represent an arbitrary matrix $R$ by its masking descriptor $b$.
If $\mathsf{LB}(R)$ holds, apply it at $h=Q(b)$ to obtain a bound
$z$.  The bad-branch conjunction $r\land e\land\neg f$ is false
for that $z$: if $r$ holds then $P(b,z)$ is a subset of $z$, so
$e$ implies $f$.  The three hull clauses consequently hold for
every $x,y$.

When $x$ is empty, $a=E(W(x,d),x)$ is false; the first clause
therefore supplies closure at empty inputs.  The second gives
closure at members of $h$.  The last gives support for each
$y\in h$, with input witness $P(b,y)$.  These are precisely the
two conditions in~\eqref{eq:e639-epsc}.\hfill$\square$

\subsection{The converse model construction still descends}

\paragraph{Upper-model theorem.}
From any countable nonempty raw-EPSC structure one can construct
a nontrivial model of the seven new equations.

\paragraph{Proof.}
Let $(M,E_M)$ be such a structure and fix $m_0\in M$.
Use the inherited free carrier $D$ of finite trees with distinct
leaves $d$ and $e_m$ for $m\in M$, and disjoint constructors
$P,A,Q$.  The leaves $e_m$ are elements of the model, not additional
symbols in the packet.  Interpret those three operations as actual
free-tree constructors.  Give all leaves weight one, and set
\begin{equation}
 w(Q(b))=w(b)+1,\quad
 w(P(b,z))=w(b)+w(z)+2,\quad
 w(A(b,z))=w(b)+w(z)+3.
 \label{eq:e639-weights}
\end{equation}

Let $\operatorname{idx}(i)$ count the consecutive outer $Q$ layers
of $i$.  A tree with another outer symbol has index zero.
Read a $P$-tree as a cons cell and a non-$P$ tail as the infinite
constant-$d$ stream.  Define $V(s,i)$ by reading coordinate
$\operatorname{idx}(i)$.  Head and Tail hold for all arguments.

Decode a class tree by syntax: each leaf denotes falsity; $P$
is NAND; $A(i,j)$ is membership between its two indexed coordinates;
and $Q(a)$ prepends a coordinate and existentially quantifies the
negation of its body.  This finite formula decoding uses no raw
value assignment.

The changed descriptor relation is now the decoded class of $b$
itself, evaluated on $P(y,P(x,b))$.  It is a finite pure-membership
formula $\widehat R_b(x,y;\bar t)$.  Every fixed parameter
$t\in\bar t$ is a proper subtree of $b$ or the default leaf $d$.
To verify this, follow each finite class occurrence and each of its
selector indices.  A referenced coordinate is a bound variable,
one of the two inputs, or a coordinate of the native stream $b$.
In the last case it is the head of a proper $P$-tail subtree of
$b$, or $d$ after the finite spine ends.  The full class syntax
may be $b$, but it does not require the raw value of $b$ merely
to inspect that syntax.  This is the crucial descent fact that
must be rechecked after dropping $V(b,d)$.

Define a surjection $\mu:D\to M$ by recursion on the positive
weights, starting with $\mu(d)=m_0$ and $\mu(e_m)=m$.
For a descriptor $b$, its proper-subtree parameters already have
raw values when its hull $Q(b)$ is assigned.  Let $R_b^M$ be the
resulting relation on $M$.

If $R_b^M$ is locally subset-bounded, choose an EPSC hull $h_b$.
Otherwise classical negation of local boundedness gives an input
$h_b$ such that for every $z\in M$ there are $n,m\in M$ with
\[
 R_b^M(h_b,n)\land E_M(m,n)\land\neg E_M(m,z).
\]
Choose such an input and set $\mu(Q(b))=h_b$.

In the unbounded case, when assigning $P(b,z)$ choose its raw
value $n$ so that $R_b^M(h_b,n)$ and
$n\not\subseteq_{E_M}\mu(z)$.  When assigning $A(b,z)$ choose
a member of this $n$ outside $\mu(z)$.
In the hull case, if $E_M(\mu(z),h_b)$, choose the raw value
$n$ of $P(b,z)$ so that $R_b^M(n,\mu(z))$, using hull support;
otherwise use $m_0$.  Give $A(b,z)$ the default value $m_0$
in the hull case.

The recursion is well founded: the raw values required for a
descriptor's parameters precede $Q(b)$; both $Q(b)$ and $z$
precede $P(b,z)$; and $P(b,z)$ precedes $A(b,z)$ by
\eqref{eq:e639-weights}.  The free constructor tags give each
tree only one such assignment.  A fixed enumeration of $M$
selects the least available witnesses, making the choices uniform.
Every requested witness exists by the selected bounded or
unbounded case.

Set $E_D(u,v)\iff E_M(\mu(u),\mu(v))$.  Surjectivity of $\mu$
and induction on formulas show that every equality-free formula
is preserved and reflected under $\mu$, with the corresponding
raw parameter values.  Native equality between trees is never
replaced by equality of raw projections.

Let $\varepsilon(s,a)$ be the truth value of the decoded finite
class in $(D,E_D)$.  For $a\ne d$, choose $W(s,a)$ to falsify
$\varepsilon(P(z,s),a)$ whenever a false instance exists, and
choose $d$ otherwise.  At $a=d$ the class is always false,
so the binder equation holds regardless of $W(s,d)$.  Use this
free slice to choose a member of the raw object represented by
$s$ whenever it is nonempty, and $d$ otherwise.  In particular,
$E_D(W(x,d),x)$ is true whenever $x$ is nonempty.

In the unbounded case the chosen raw values of
$Q(b),P(b,z),A(b,z)$ make every bad branch true.  In the hull
case EPSC and the selected member and support witnesses make
all three hull clauses true.  Thus the last equation holds for
all tuples.  This verifies the changed closure relation itself,
not an unmodified predecessor relation.

Finally transport a countably infinite Boolean algebra, such as
the finite/cofinite algebra of subsets of $\mathbb N$, onto $D$
with top at its distinguished element $d$.  Interpret $N$ as
NAND and $p_E,p_\varepsilon$ by Boolean top or bottom according
to $E_D,\varepsilon$.  Retain all native functions.  Every equation
holds, and $Q(d)\ne d$ because the constructors are free.
This supplies the required nontrivial model.\hfill$\square$

\subsection{Noncollapse, formalization base, and the new calibration}

The native Head and Tail equations alone show that $Q(d)=d$
collapses every model to a singleton.  Indeed,
\[
 x=V(P(x,s),d)=V(P(x,s),Q(d))=V(s,d)
\]
for every $x,s$.  Conversely, every singleton model satisfies
the query.  Sh1 likewise shows that $N(d,d)=d$ holds exactly
in singleton models, since it asserts $\neg d=d$ in the
Boolean reduct.

For completeness, the same argument has a useful general form.
In any nontrivial Head/Tail model the selector values
$d,Q(d),Q^2(d),\ldots$ are pairwise distinct.  If two indices
$m<n$ coincide, evaluate a finite stack whose coordinates
$m,n$ are arbitrary elements $x,y$.  The coordinate lemma
gives $x=y$, contradicting nontriviality.  In particular no
finite nontrivial structure can model even this native fragment.
Finite regression checks of encodings must not be treated as
finite consistency evidence.

Let $\mathcal E_{639}$ denote the seven equations.  Standard
ground equational completeness gives
\[
 \mathcal E_{639}\nvdash Q(d)=d
 \quad\Longleftrightarrow\quad
 \mathcal E_{639}\text{ has a nontrivial model}.
\]
One can use its elementary proof: take the least congruence
on ground terms containing every ground axiom instance.  If
the endpoint classes are distinct, the quotient is a countable
separating model.  The constant $d$ makes its domain nonempty.
Conversely equational soundness preserves any separating model.
Each rewrite certificate must assign values to the union of
variables on both sides, including variables absent from the
side currently being matched.

The two new model theorems prove in ordinary ZFC
\begin{equation}
 \mathcal E_{639}\nvdash Q(d)=d
 \quad\Longleftrightarrow\quad
 \text{raw EPSC has a model}
 \quad\Longleftrightarrow\quad
 \mathcal E_{663}\nvdash Q(d)=d.
 \label{eq:e639-relative}
\end{equation}
For an arbitrary raw-EPSC model one may first take a countable
elementary submodel in its countable language.  First-order
completeness supplies the corresponding consistency statement.
The raw-EPSC/ZFC comparison is inherited from the expressly pinned
ProveIt Sources 08 and 14 and the E692 compatibility proof, not
reproved by the finite tests here.\footnote{The source dossier is
\url{https://github.com/VladimirReshetnikov/ProveIt/blob/ce192939372ac7944fe2a0208c68a96d49fb5e3b/SetTheory/Cardinals/docs/reports/set-theory/zfc-closure-kernels/intake-146-source-14.md}.
The retained E663/E667 packet is from the user-specified ProveIt
commit \texttt{0128d662fd}; its reference directory includes
\texttt{e692\_compatibility.tex} and the raw-EPSC source snapshot.}
With that inherited calibration,
\begin{equation}
 \mathrm{ZFC}\vdash
 \left(\mathcal E_{639}\nvdash Q(d)=d
              \ \longleftrightarrow\ \operatorname{Con}(\mathrm{ZFC})\right).
 \label{eq:e639-zfc}
\end{equation}
Equivalently, the finite first-order theory
$\mathcal E_{639}+\{Q(d)\ne d\}$ has this consistency strength.

All constructions used in the new proof are set constructions,
recursions on finite syntax or positive integer weights, ordinary
Boolean algebra, and first-order soundness and completeness.
They are available in ZFC.  No model of ZFC is asserted to exist
inside ZFC without the consistency premise, no universe axiom
is added, and no general soundness of a proof assistant is used.
The result is an ordinary mathematical proof formalizable in ZFC.
The artifact does not claim a newly serialized Lean/Rocq endpoint
or a represented PA proof-code transformer.

\subsection{Exact packet accounting and reproducible checks}

The equation-node counts are now
\[
                 14,6,9,12,12,17,74,
\]
totaling 144.  The seven equations contain four distinct local
variables at most; nine operation symbols plus those four
variables still require four-bit term opcodes.  The complete
packet count is
\[
 \underbrace{8+3\cdot9+8}_{43\text{ header bits}}
 +\underbrace{7+1+4\cdot144}_{584\text{ equation bits}}
 +\underbrace{4\cdot3}_{12\text{ endpoint bits}}
 =\boxed{639\text{ bits}}.
\]
It occupies 80 transport bytes with one trailing zero padding
bit, excluded from the logical packet.  Its exact bit-character
SHA-256, excluding the final newline, is
\begin{center}\scriptsize\ttfamily
15c8bde20963ac71ff9015556c52e88ac3158e9d0b585f056cd11f018ed006a8.
\end{center}

The self-contained \texttt{build\_e639.py} imports no ancestor
Python code.  It reconstructs the exact pinned E663 bits, checks
that the signature, first six equations, and endpoints are unchanged,
and emits the new full AST, seven expanded equations, bit string,
packed bytes, bit length, and verification receipt.  Its independent
decoder round-trips the packet and rejects all 639 strict truncations
and three malformed controls.

The new masking identities are checked on 1,020 Boolean assignments
through seven native parameters, and 4,096 assignments in the
eight-element powerset Boolean algebra.  Another 1,155 exact native
coordinate checks cover up to 32 parameters.  All 512 assignments
to the inherited Boolean closure kernel pass its cofactor and
truth-forcing checks.  These are transcription and regression checks;
the all-model proofs above establish the unbounded conclusions.

This construction places some additional syntax in possible proof
certificates: representing a matrix uses the duplicated parameter
spine.  That syntax is generated by a fair proof search, rather
than stored in the literal 639-bit theory.  Whether this change
reduces a complete binary machine is an implementation question.
The next useful experiment is a complete ground-rewrite interpreter
that charges the packet writer, substitution, context traversal,
and fair enumeration together.
